Ze heette peettante,
maar ook roddeltante.

Haar moeder was een hoer.
Haar vader was een boer.

De ouwe theetante,
zo oud als een tante,

was zijn echtelijk kind.
Moeder was zijn kleinkind,

maar ook zijn nageslacht.
Hij had Tante verkracht

en slordig bezwangerd.
De schuldige bangerd

eiste dat Tante loog
en zo de kerk bedroog.

Ook de prostituee
deed aan de leugen mee,

nam moeder als kind aan
en liet haar man begaan.

Ze nam de oplichting,
daarna mee háár graf in,

als een Tante Meier,
een dode temeier.

Het familiegeheim
verstikte in zijn slijm.

Tante zag denken als
iets beperkts en vaak vals.

Ze meed breinalgebra.
“Voel eerst eens twee keer na

voordat je reageert!”,
had ze moeder geleerd.

“Ik snap dat je ouders
beschouwt als domhouders,

maar het is wel slim nog
om ze te hebben toch!

Pa gaf zijn pijnstokje
door, als een klein vlokje,

aan ons als kinderen.
Dit te verminderen

is onze levenstaak!
De oorsprong is bijzaak.”

Moeder huilde hardop.
Tante zei toen daarop:

“Je scheidt meer van jezelf
dan van je wederhelft.
Je hebt je man nodig.
Lust is overbodig!”

Tante was geen Hegel,
maar had als leefregel

“Maak het niet te moeilijk!
God is te aanzienlijk

om door ons te vatten!”
Zo kon ze soebatten:

“Onze integriteit
leidt tot waarachtigheid

en vraagt liefdevol zijn
voor je zelfbewustzijn.”

Oom werd vaak afgekat.
“Als tante ballen had

dan was ze meer een oom.”
Oom was bangig en sloom.

Oom was mentaal zachter.
Oom was een halfslachter.

Zijn gebaren van macht
misten telkens de kracht

van uiteindelijk doen.
Emmy vond hem een oen,
stevig onder de plak.
Een schichtige goedzak.

In plaats van afdwingen
kon hij slootjespringen

als iemand hem dreigde.
Want oom verstijfde.

Als tante werd bespot
dan werd Emmy beknot

als ze voor haar opkwam.
Wat slechts in hem op kwam

was de vrees voor een strijd
en eigen veiligheid

in plaats van tante’s eer.
Hij trachtte ooit een keer

wel moedig te kiezen
voor zijn Eurydice,

maar trok ze laf terug
na een tik op zijn rug.

De slappe antiheld
werd als Orpheus geveld.

Zijn zoon was het herstel
van zijn kommer en kwel.

Hij was zijn enige.
Schoon hij het menige

keren had geprobeerd,
werd hij telkens geweerd

door haar ferme tante.
Naar de bijdehante

was het genoeg geweest.
Haar oom was te bevreesd

voor schande op zijn werk
om dit feit aan hun kerk

te melden als verzuim.
Dicht ging tante’s zuurpruim.

Middelpuntvliedende
klachten of ziedende

uitspattingen ten spijt.
Hij was zijn lusthol kwijt.

Tante negeerde totaal
haar ooms gebarenpaal.

Hij had niet gemekkerd,
maar zijn ‘slaap lief lekkerd’

was niet meer te horen.
Hij mocht haar niet storen

als Tante naar bed ging.
Ze zocht afzondering.

Oom kon die gehechtheid
niet in zijn geloof kwijt,

zich open uitlaten
of leren loslaten.

De Aparigraha
bekend uit de tantra,

als het niet-vergaren,
kon hij niet verklaren

als het ging om de daad.
Want het scheppingsmandaat

was juist Gods instructie
en dit was obstructie.

Haar oom zijn frustraties
over zijn prestaties

moest zijn zoon verzetten.
Zijn falen rechtzetten.

Hij maakte ooms droom waar.
En werd van misdienaar

zelfs een hofkapelaan.
Vader hoorde hem aan

staand naast de prins-bisschop
als zijn misschaduw op

het balkon van de vorst.
Al kloppend op zijn borst

verkreeg ook hij status
in zijn catechismus.

Hij was door te slikken
en altijd ja-knikken

wel wat opgeklommen,
maar wist te verstommen

bij het geestig talent
dat zijn zoon had verkend.

Geroemd om zijn vroomheid,
dat in zijn directheid

zelfs de bisschop versloeg.
Toen de kapelaan vroeg

“hoe vaak heeft u per jaar
als gehoorzaam echtpaar

uw echtelijke plicht
lichamelijk verricht?”

stond de bisschop eerst paf.
De kapelaan telde af.

Tot zijn verwondering
stak na wat huivering

iemand uit het publiek
zijn hand op als repliek

bij het getal duizend.
Het per dag uitpluizend

ontsloeg de menigte,
die zich verenigde

in kleinerend gelach
of afgunstig ontzag.

De eerstvolgende hand
kwam veel later tot stand.

De bisschop begreep toen
deze manier van doen.

“Meer communicatie
door participatie!

Dat heb je slim gedaan”,
hoorde de kapelaan

hem zacht toefluisteren.
Het volk bleef luisteren

naar de wipgetallen.
Het was welgevallen

de onderlinge strijd
tussen trots en waarheid.

Pas bij tien werd het sneu.
Toch was men het niet beu

om elkaar te jennen.
Niemand liet zich kennen
tot het laatste getal.
Wat wilde het geval

drie handen staken uit.
De één had net zijn bruid.

Dus dat verklaarde veel.
Weg ging het gekrakeel

door het medelijden
met het schijnbaar lijden

van twee uit de partij.
Eentje keek juist heel blij

en begon te juichen
en dank te betuigen.

“Vandaag is juist de dag
dat ik er weer op mag!”,

jubelde hij verrukt.
De ander keek bedrukt.

“Ik ook”, zei die laatste
toen hij de bal kaatste

op, krachtens zijn grimas,
waarom hij treurig was?

Hij wees toen op zijn vrouw.
Men begreep hem al gauw.

Tante ontving Emmy
pas na een absentie

van vele jaren weer.
Het was steeds de landheer

die zo’n bezoek verbood.
Dus kwam Tante uit nood

altijd bij moeder langs.
Haar tante had onlangs

een flinke griep gevat,
die haar afgemat had.

Emmy mocht haar bijstaan
en ditmaal alleen gaan.

Oom was in het klooster
volgens zijn dienstrooster

en sprak daar haar neef toe.
Emmy ging er naar toe.

Zijn zoon werd gewaarschuwd.
“Je bent wel ongehuwd,

maar hou je eer intact.
Menig man raakt verzwakt

door een publieke vrouw.
Blijf je celibaat trouw.

Als een vrouw je meevraagt,
juist als ze je behaagt,

wijs haar dan kordaat af.
Dat is haar juiste straf.”

Haar neef liep langs Emmy,
maar gaf haar geen notie.

Tot ze vroeg: "ga je mee?”
Hij zei meteen: “Nee, nee!

Ik wil geen sex met je!”
Kon het flauwe tekstje,

toen hij haar herkende
en zij zich afwende,

zijn excuus vertellen
en de band herstellen?
